\documentclass[12pt,a4paper]{article}
\usepackage[utf8]{inputenc}
\usepackage[T5]{fontenc}
\usepackage[vietnamese]{babel}
\usepackage[margin=1.5cm,top=1.4cm,bottom=1.4cm]{geometry}
\usepackage{amsmath,amssymb,array,tabularx,enumitem,fancyhdr,tikz,xcolor}
\usetikzlibrary{calc,angles,quotes,arrows.meta,babel,patterns}
\definecolor{ink}{RGB}{30,55,70}

\newcommand{\headerQuote}{%
  \ifcase\value{page}\relax
    \or Math gives us hope that every problem has a solution% Trang 1
    \or Math + You = Champion% Trang 2
    \or Where Frustration Becomes Inspiration% Trang 3
    \or Turning Confusion into Clarity% Trang 4
    \or Don't be an X, be a Y of mathematics% Trang 5
    \or Don't get mad. Learn to add% Trang 6
    \or Math gives us hope that every problem has a solution% Trang 7
    \or Math + You = Champion% Trang 8
    \or Where Frustration Becomes Inspiration% Trang 9
    \or Turning Confusion into Clarity% Trang 10
    \or Don't be an X, be a Y of mathematics% Trang 11
    \or Don't get mad. Learn to add% Trang 12
    \else Math gives us hope that every problem has a solution%
  \fi
}
\pagestyle{fancy}\fancyhf{}
\lhead{\small\textit{\headerQuote}}
\rhead{\small\textit{Người biên soạn: Dương Gia Bảo}}
\rfoot{\small Trang \thepage}
\renewcommand{\headrulewidth}{0.4pt}
\setlength{\headheight}{18pt}
\setlength{\parindent}{0pt}
\setlength{\parskip}{1.5pt}
\setlength{\tabcolsep}{4pt}
\renewcommand{\arraystretch}{1.12}
\newcommand{\parttitle}[1]{\par\vspace{2pt}{\large\bfseries\color{ink}#1}\par\vspace{1pt}}
\newcommand{\four}[4]{\par\noindent\begin{tabularx}{\linewidth}{@{}*{4}{>{\raggedright\arraybackslash}X}@{}}\textbf{A.} #1&\textbf{B.} #2&\textbf{C.} #3&\textbf{D.} #4\end{tabularx}}
\newcommand{\two}[4]{\par\noindent\begin{tabularx}{\linewidth}{@{}*{2}{>{\raggedright\arraybackslash}X}@{}}\textbf{A.} #1&\textbf{B.} #2\\[2pt]\textbf{C.} #3&\textbf{D.} #4\end{tabularx}}
\newcommand{\answerbox}{\hfill\textit{Đáp số:}\ \fbox{\rule{0pt}{12pt}\hspace{2.2cm}}}
\setlist[enumerate,1]{label=\textbf{Câu \arabic*.},leftmargin=1.4cm,labelsep=.15cm,itemsep=1.5pt,topsep=1.0pt,parsep=0pt}
\setlist[enumerate,2]{label=\alph*),leftmargin=.6cm,itemsep=1.0pt,topsep=1.0pt,parsep=0pt}

\begin{document}
\noindent
\begin{minipage}[c]{0.38\textwidth}
    \centering
    \textbf{LÀM TOÁN ĐỂ HẠNH PHÚC}\\[-2pt]
    \rule{3.2cm}{0.6pt}\\[4pt]
    \textbf{ĐỀ CHÍNH THỨC}
\end{minipage}\hfill
\begin{minipage}[c]{0.35\textwidth}
    \centering
    \textbf{ĐỀ ÔN THI GK1}\\
    \textbf{NĂM HỌC 2026 -- 2027}\\[2pt]
    \textbf{ĐỀ SỐ 09}\\
    \textbf{Môn: TOÁN -- Khối lớp 10}\\
    \textit{\small Thời gian làm bài: 90 phút\\(không kể thời gian phát đề)}
\end{minipage}

\vspace{0.2cm}
\noindent
\begin{minipage}[b]{0.78\textwidth}
    \textbf{Họ và tên:}\dotfill\textbf{SBD:}\makebox[3.5cm]{\dotfill}
\end{minipage}\hfill
\begin{minipage}[b]{0.20\textwidth}
    \hfill
    \setlength{\fboxrule}{0.8pt}%
    \fbox{\begin{minipage}{2.3cm}
        \centering
        \textbf{\small Mã đề thi}\\
        \textbf{0009}
    \end{minipage}}
\end{minipage}
\vspace{0.2cm}

\parttitle{PHẦN I. Trắc nghiệm nhiều lựa chọn}
\textit{Thí sinh trả lời từ câu 1 đến câu 12. Mỗi câu hỏi chỉ chọn một phương án đúng duy nhất.}
\begin{enumerate}
\item Trong các mệnh đề sau đây, mệnh đề nào là mệnh đề \textbf{đúng}?
\four{Số $9$ là một số nguyên tố.}{$\pi < 3$.}{Tổng ba góc trong một tam giác bằng $180^\circ$.}{$\sqrt{2}$ là một số hữu tỉ.}

\item Cho tập hợp $A = \{1; 2; 3\}$. Khẳng định nào sau đây là \textbf{sai}?
\four{$1 \in A$.}{$\{2\} \in A$.}{$\{1; 2\} \subset A$.}{$\varnothing \subset A$.}

\item Trong các câu sau, câu nào là một mệnh đề toán học?
\four{Hôm nay trời đẹp quá!}{$2 + 3 = 7$.}{Bạn có thích học môn Toán không?}{Hãy làm bài thật cẩn thận nhé!}

\item Tập hợp gồm tất cả các số thực $x$ thỏa mãn $2 \le x < 6$ được viết dưới dạng khoảng, đoạn, nửa khoảng là
\four{$[2; 6)$.}{$(2; 6]$.}{$[2; 6]$.}{$(2; 6)$.}

\item Cặp số $(x; y)$ nào sau đây là một nghiệm của hệ bất phương trình $\begin{cases} 2x - y \ge 1 \\ x + y \le 4 \end{cases}$?
\four{$(0; 0)$.}{$(1; 3)$.}{$(2; 1)$.}{$(0; 2)$.}

\item Điểm nào sau đây thuộc miền nghiệm của bất phương trình $x - 3y + 3 > 0$?
\four{$M(0; 2)$.}{$N(1; 2)$.}{$P(-4; 0)$.}{$Q(1; 0)$.}

\item Miền nghiệm của bất phương trình $2x + y - 4 \le 0$ là nửa mặt phẳng (kể cả bờ $d: 2x + y - 4 = 0$) chứa điểm nào sau đây?
\four{$A(3; 1)$.}{$O(0; 0)$.}{$B(2; 2)$.}{$C(1; 3)$.}

\item Đường thẳng bờ của miền nghiệm của bất phương trình $3x - 2y + 6 \ge 0$ có phương trình là
\four{$2x - 3y + 6 = 0$.}{$3x + 2y - 6 = 0$.}{$3x - 2y = 0$.}{$3x - 2y + 6 = 0$.}

\item Cho góc $\alpha$ là một góc tù ($90^\circ < \alpha < 180^\circ$). Khẳng định nào sau đây là \textbf{đúng}?
\four{$\sin \alpha > 0$.}{$\cos \alpha > 0$.}{$\tan \alpha > 0$.}{$\cot \alpha > 0$.}

\item Cho tam giác $ABC$ có cạnh $b = 6$, $c = 8$ và góc xen giữa $\widehat{A} = 60^\circ$. Độ dài cạnh $a = BC$ bằng
\four{$2\sqrt{13}$.}{$10$.}{$2\sqrt{37}$.}{$\sqrt{52}$.}

\item Giá trị của biểu thức $P = \sin^2 40^\circ + \sin^2 50^\circ + \tan 30^\circ \cdot \cot 30^\circ$ bằng
\four{$1$.}{$3$.}{$2$.}{$0$.}

\item Giá trị của biểu thức lượng giác $\sin 150^\circ$ bằng
\four{$-\dfrac{1}{2}$.}{$\dfrac{\sqrt{3}}{2}$.}{$-\dfrac{\sqrt{3}}{2}$.}{$\dfrac{1}{2}$.}
\end{enumerate}

\parttitle{PHẦN II. Trắc nghiệm đúng/sai}
\textit{Thí sinh trả lời từ câu 13 đến câu 14. Trong mỗi ý a), b), c), d) ở mỗi câu, thí sinh chọn đúng hoặc sai.}
\begin{enumerate}[resume]
\item Cho hệ bất phương trình bậc nhất hai ẩn $\begin{cases} x \ge 0,\; y \ge 0 \\ x + y \le 4 \\ x + 2y \le 6 \end{cases}$. Biết miền nghiệm của hệ là miền tứ giác $OABC$ trên mặt phẳng tọa độ $Oxy$.
\begin{enumerate}
    \item Điểm $M(1; 1)$ thuộc miền nghiệm của hệ bất phương trình đã cho.
    \item Tọa độ đỉnh $A$ nằm trên trục hoành $Ox$ là $A(4; 0)$.
    \item Tọa độ đỉnh $B$ là giao điểm của hai đường thẳng $x + y = 4$ và $x + 2y = 6$ là $B(2; 2)$.
    \item Biểu thức $F(x, y) = 2x + 3y$ đạt giá trị lớn nhất trên miền nghiệm bằng $12$.
\end{enumerate}

\item Cho tam giác $ABC$ có độ dài ba cạnh lần lượt là $a = BC = 7$, $b = CA = 5$, $c = AB = 6$.
\begin{enumerate}
    \item Nửa chu vi của tam giác $ABC$ là $p = 9$.
    \item Diện tích tam giác $ABC$ tính theo công thức Heron bằng $6\sqrt{6}$.
    \item Bán kính đường tròn nội tiếp của tam giác $ABC$ là $r = \dfrac{2\sqrt{6}}{3}$.
    \item Bán kính đường tròn ngoại tiếp của tam giác $ABC$ là $R = \dfrac{35\sqrt{6}}{12}$.
\end{enumerate}
\end{enumerate}

\parttitle{PHẦN III. Trắc nghiệm trả lời ngắn}
\textit{Thí sinh trả lời từ câu 15 đến câu 18. Viết kết quả vào ô quy định.}
\begin{enumerate}[resume]
\item Lớp 10A có $40$ học sinh, trong đó có $25$ bạn thích môn Toán, $20$ bạn thích môn Ngữ văn, và $5$ bạn không thích cả hai môn này. Hỏi có bao nhiêu học sinh thích cả hai môn Toán và Ngữ văn?
\answerbox

\item Cho hệ bất phương trình bậc nhất hai ẩn $\begin{cases} x \ge 0,\; y \ge 0 \\ x + y \le 5 \\ 3x + y \le 9 \end{cases}$.
Tìm giá trị lớn nhất của biểu thức $F(x, y) = 4x + 3y$ trên miền nghiệm của hệ.
\answerbox

\item Cho góc nhọn $\alpha$ thỏa mãn $0^\circ < \alpha < 90^\circ$ và $\sin \alpha = \dfrac{5}{13}$. Tính giá trị của biểu thức $T = 13\cos \alpha + 12\tan \alpha$.
\answerbox

\item \begin{minipage}[t]{0.62\linewidth}
Để xác định khoảng cách giữa hai vị trí $A$ và $B$ bị ngăn cách bởi một hồ nước sâu, một người đứng tại vị trí $C$ đo được khoảng cách $CA = 60\text{ m}$, $CB = 100\text{ m}$ và góc $\widehat{ACB} = 120^\circ$. Hỏi khoảng cách giữa hai vị trí $A$ và $B$ bằng bao nhiêu mét? (\textit{Làm tròn kết quả đến hàng đơn vị}).
\vspace{6pt}

\answerbox
\end{minipage}\hfill
\begin{minipage}[t]{0.35\linewidth}
\centering
\begin{tikzpicture}[scale=0.75, every node/.style={scale=0.85}]
    \coordinate (C) at (0,0);
    \coordinate (A) at (-2.0, 1.2);
    \coordinate (B) at (3.5, 0);

    % Lake representation
    \draw[blue!40, fill=blue!8, dashed] (0.8,0.7) ellipse (1.6cm and 0.7cm);
    \node[blue!70!black, font=\scriptsize\itshape] at (0.8,0.7) {Hồ nước};

    \draw[thick, ink] (C) node[below] {$C$} -- (A) node[above left] {$A$};
    \draw[thick, ink] (C) -- (B) node[below right] {$B$};
    \draw[dashed, red!80!black, line width=1.1pt] (A) -- (B);

    \filldraw[black] (C) circle (1.8pt);
    \filldraw[black] (A) circle (1.8pt);
    \filldraw[black] (B) circle (1.8pt);

    \node at (-1.3,0.3) {\small $60\text{ m}$};
    \node at (1.8,-0.35) {\small $100\text{ m}$};
    \node at (0.3,0.3) {\small $120^\circ$};
    \node[above, red!80!black] at (0.8,1.4) {$d = ?$};
\end{tikzpicture}
\end{minipage}
\end{enumerate}

\parttitle{PHẦN IV. Tự luận}
\textit{Thí sinh trình bày lời giải chi tiết từ câu 19 đến câu 22.}
\begin{enumerate}[resume]
\item \textit{(1,0 điểm)} Cho hai tập hợp khác rỗng $A = [m - 3; m + 2]$ và $B = [-1; 5)$ với $m$ là tham số thực. Tìm tất cả các giá trị của $m$ để tập hợp $A$ là tập con của tập hợp $B$ ($A \subset B$).

\item \textit{(0,75 điểm)} Cho tam giác $ABC$ có độ dài ba cạnh lần lượt là $a = 13$, $b = 14$ và $c = 15$.
\begin{enumerate}
    \item Tính diện tích tam giác $ABC$.
    \item Tính độ dài đường cao $h_a$ kẻ từ đỉnh $A$ của tam giác $ABC$.
\end{enumerate}

\item \begin{minipage}[t]{0.58\linewidth}
\textit{(0,75 điểm)} Để đo chiều cao của một cột phát sóng viễn thông $CD$ dựng thẳng đứng trên mặt đất phẳng ($D$ là chân cột, $C$ là đỉnh cột), người ta chọn hai điểm quan sát $A$ và $B$ trên mặt đất sao cho ba điểm $A, B, D$ thẳng hàng theo thứ tự đó, với khoảng cách $AB = 30\text{ m}$. Từ $A$, góc nâng nhìn lên đỉnh cột là $\widehat{CAD} = 35^\circ$. Từ $B$, góc nâng nhìn lên đỉnh cột là $\widehat{CBD} = 55^\circ$. Tính chiều cao $h = CD$ của cột phát sóng (\textit{làm tròn kết quả đến một chữ số thập phân}).
\end{minipage}\hfill
\begin{minipage}[t]{0.38\linewidth}
\centering
\begin{tikzpicture}[scale=0.82, every node/.style={scale=0.85}]
    \coordinate (A) at (0,0);
    \coordinate (B) at (2.2,0);
    \coordinate (D) at (4.5,0);
    \coordinate (C) at (4.5,3.0);

    \draw[thick] (-0.4,0) -- (5.0,0);
    \draw[thick, ink] (A) node[below] {$A$} -- (B) node[below] {$B$} -- (D) node[below] {$D$};
    \draw[line width=1.4pt, blue!80!black] (D) -- (C);
    \draw[dashed, teal!80!black, thick] (A) -- (C) -- (B);

    \filldraw[black] (A) circle (1.5pt);
    \filldraw[black] (B) circle (1.5pt);
    \filldraw[black] (D) circle (1.5pt);
    \filldraw[red!80!black] (C) circle (2.2pt) node[above] {$C$};

    \draw (4.25,0) -- (4.25,0.25) -- (4.5,0.25);
    \node at (1.1,-0.35) {\small $30\text{ m}$};
    \node at (0.8,0.25) {\small $35^\circ$};
    \node at (2.8,0.3) {\small $55^\circ$};
    \node[right, blue!80!black] at (4.5, 1.5) {$h$};
\end{tikzpicture}
\end{minipage}

\item \textit{(0,5 điểm)} Một xưởng sản xuất bàn ghế nội thất sản xuất hai loại mặt hàng: bàn gỗ (loại I) và ghế gỗ (loại II).
\begin{itemize}[leftmargin=0.5cm]
    \item Để sản xuất một chiếc bàn gỗ cần $2\text{ m}^2$ gỗ và $3$ giờ công lao động.
    \item Để sản xuất một chiếc ghế gỗ cần $1\text{ m}^2$ gỗ và $2$ giờ công lao động.
\end{itemize}
Xưởng có nguồn nguyên liệu dự trữ là $80\text{ m}^2$ gỗ và tối đa $130$ giờ công lao động trong một tuần. Mỗi chiếc bàn bán ra đem lại lợi nhuận $400\,000$ đồng, mỗi chiếc ghế đem lại lợi nhuận $250\,000$ đồng. Hỏi xưởng cần sản xuất bao nhiêu chiếc bàn và bao nhiêu chiếc ghế trong một tuần để tổng lợi nhuận thu được là lớn nhất?
\end{enumerate}

\newpage
% =========================================================================
% BẢNG ĐÁP ÁN VÀ LỜI GIẢI CHI TIẾT
% =========================================================================
\parttitle{BẢNG ĐÁP ÁN TỔNG HỢP}

\vspace{0.2cm}
\textbf{1. Phần I: Trắc nghiệm nhiều lựa chọn}
\begin{center}
{\setlength{\tabcolsep}{2.5pt}\small
\begin{tabular}{|c|c|c|c|c|c|c|c|c|c|c|c|}
\hline
\textbf{Câu 1} & \textbf{Câu 2} & \textbf{Câu 3} & \textbf{Câu 4} & \textbf{Câu 5} & \textbf{Câu 6} & \textbf{Câu 7} & \textbf{Câu 8} & \textbf{Câu 9} & \textbf{Câu 10} & \textbf{Câu 11} & \textbf{Câu 12} \\
\hline
\textbf{C} & \textbf{B} & \textbf{B} & \textbf{A} & \textbf{C} & \textbf{D} & \textbf{B} & \textbf{D} & \textbf{A} & \textbf{A} & \textbf{C} & \textbf{D} \\
\hline
\end{tabular}}
\end{center}

\vspace{0.2cm}
\textbf{2. Phần II: Trắc nghiệm đúng/sai}
\begin{center}
\begin{tabular}{|c|c|c|c|c|}
\hline
\textbf{Câu hỏi} & \textbf{Ý a)} & \textbf{Ý b)} & \textbf{Ý c)} & \textbf{Ý d)} \\
\hline
\textbf{Câu 13} & \textbf{Đúng} & \textbf{Đúng} & \textbf{Đúng} & \textbf{Sai} \\
\hline
\textbf{Câu 14} & \textbf{Đúng} & \textbf{Đúng} & \textbf{Đúng} & \textbf{Sai} \\
\hline
\end{tabular}
\end{center}

\vspace{0.2cm}
\textbf{3. Phần III: Trắc nghiệm trả lời ngắn}
\begin{center}
\begin{tabular}{|c|c|c|c|}
\hline
\textbf{Câu 15} & \textbf{Câu 16} & \textbf{Câu 17} & \textbf{Câu 18} \\
\hline
\textbf{10} & \textbf{17} & \textbf{17} & \textbf{140} \\
\hline
\end{tabular}
\end{center}

\vspace{0.2cm}
\textbf{4. Phần IV: Tự luận}
\begin{center}
\begin{tabular}{|c|l|}
\hline
\textbf{Câu hỏi} & \textbf{Kết quả tóm tắt} \\
\hline
\textbf{Câu 19} & $m \in [2; 3)$ hay $2 \le m < 3$. \\
\hline
\textbf{Câu 20} & a) $S = 84$;\quad b) $h_a = \dfrac{168}{13} \approx 12{,}9\text{ (cm hoặc đơn vị độ dài)}$. \\
\hline
\textbf{Câu 21} & Chiều cao cột phát sóng $h \approx 41{,}2\text{ m}$. \\
\hline
\textbf{Câu 22} & Sản xuất \textbf{30 chiếc bàn} và \textbf{20 chiếc ghế}; lợi nhuận lớn nhất là $17\,000\,000$ đồng. \\
\hline
\end{tabular}
\end{center}

\vspace{0.4cm}
% =========================================================================
% HÌNH VẼ MIỀN NGHIỆM ĐỘC LẬP NGOÀI BẢNG CHO CÂU 16 VÀ CÂU 22
% =========================================================================
\parttitle{HÌNH VẼ BIỂU DIỄN MIỀN NGHIỆM (HỆ BẤT PHƯƠNG TRÌNH)}

\begin{minipage}[t]{0.48\textwidth}
\centering
\textbf{Hình 1: Miền nghiệm Câu 16}\\[6pt]
\begin{tikzpicture}[scale=0.8, >=Stealth]
    % Trục toạ độ
    \draw[line width=1.2pt, ->] (-0.8,0) -- (4.5,0) node[right, font=\normalsize\bfseries] {$x$};
    \draw[line width=1.2pt, ->] (0,-0.8) -- (0,6.0) node[above, font=\normalsize\bfseries] {$y$};
    \node[below left, font=\normalsize\bfseries] at (0,0) {$O$};

    % Đa giác nghiệm O(0;0), A(3;0), B(2;3), C(0;5)
    \fill[fill=ink!18] (0,0) -- (3,0) -- (2,3) -- (0,5) -- cycle;
    \draw[line width=1.8pt, blue!70!black] (0,0) -- (3,0) -- (2,3) -- (0,5) -- cycle;

    % d1: x + y = 5 (cắt (5;0) và (0;5))
    \draw[line width=1.1pt, dashed, red!80!black] (-0.3, 5.3) -- (4.3, 0.7) node[above right, font=\small] {$d_1: x+y=5$};
    % d2: 3x + y = 9 (cắt (3;0) và (0;9))
    \draw[line width=1.1pt, dashed, teal!80!black] (-0.2, 7.5) -- (3.3, -0.6) node[below right, font=\small] {$d_2: 3x+y=9$};

    % Đỉnh
    \filldraw[black] (0,0) circle (3pt);
    \filldraw[black] (3,0) circle (3.5pt) node[below, font=\normalsize\bfseries] {$A(3;0)$};
    \filldraw[black] (2,3) circle (3.5pt) node[above right, font=\normalsize\bfseries] {$B(2;3)$};
    \filldraw[black] (0,5) circle (3.5pt) node[left, font=\normalsize\bfseries] {$C(0;5)$};

    % Dóng tọa độ B
    \draw[dotted, thick] (2,0) node[below, font=\small] {$2$} -- (2,3) -- (0,3) node[left, font=\small] {$3$};
\end{tikzpicture}
\end{minipage}\hfill
\begin{minipage}[t]{0.48\textwidth}
\centering
\textbf{Hình 2: Miền nghiệm Câu 22 (Quy hoạch tuyến tính)}\\[6pt]
\begin{tikzpicture}[x=0.06cm, y=0.048cm, >=Stealth]
    % Trục toạ độ
    \draw[line width=1.2pt, ->] (-8,0) -- (55,0) node[right, font=\normalsize\bfseries] {$x$};
    \draw[line width=1.2pt, ->] (0,-8) -- (0,75) node[above, font=\normalsize\bfseries] {$y$};
    \node[below left, font=\normalsize\bfseries] at (0,0) {$O$};

    % Đa giác nghiệm O(0;0), A(40;0), B(30;20), C(0;65)
    \fill[fill=ink!18] (0,0) -- (40,0) -- (30,20) -- (0,65) -- cycle;
    \draw[line width=1.8pt, blue!70!black] (0,0) -- (40,0) -- (30,20) -- (0,65) -- cycle;

    % d1: 2x + y = 80 (cắt (40;0) và (0;80))
    \draw[line width=1.1pt, dashed, red!80!black] (-3, 76) -- (44, -8) node[below right, font=\small] {$d_1: 2x+y=80$};
    % d2: 3x + 2y = 130 (cắt (43.3;0) và (0;65))
    \draw[line width=1.1pt, dashed, teal!80!black] (-4, 71) -- (48, -7) node[above right, font=\small] {$d_2: 3x+2y=130$};

    % Đỉnh
    \filldraw[black] (0,0) circle (3pt);
    \filldraw[black] (40,0) circle (3.5pt) node[below, font=\normalsize\bfseries] {$A(40;0)$};
    \filldraw[black] (30,20) circle (3.5pt) node[above right, font=\normalsize\bfseries] {$B(30;20)$};
    \filldraw[black] (0,65) circle (3.5pt) node[left, font=\normalsize\bfseries] {$C(0;65)$};

    % Dóng tọa độ B
    \draw[dotted, thick] (30,0) node[below, font=\small] {$30$} -- (30,20) -- (0,20) node[left, font=\small] {$20$};
\end{tikzpicture}
\end{minipage}

\newpage
% =========================================================================
% HƯỚNG DẪN GIẢI CHI TIẾT
% =========================================================================
\parttitle{HƯỚNG DẪN GIẢI CHI TIẾT}

\textbf{PHẦN I. Trắc nghiệm nhiều lựa chọn}
\begin{itemize}[leftmargin=0.6cm, itemsep=3pt]
    \item \textbf{Câu 1:} Mệnh đề toán học là khẳng định mang tính đúng hoặc sai về một chân lý toán học. Tổng ba góc trong tam giác luôn bằng $180^\circ$ là mệnh đề đúng. Chọn \textbf{C}.
    \item \textbf{Câu 2:} Phần tử $2 \in A$, do đó $\{2\} \subset A$. Kí hiệu $\{2\} \in A$ là sai. Chọn \textbf{B}.
    \item \textbf{Câu 3:} $2 + 3 = 7$ là một khẳng định toán học (mệnh đề sai). Chọn \textbf{B}.
    \item \textbf{Câu 4:} Tập hợp $\{x \in \mathbb{R} \mid 2 \le x < 6\}$ viết dưới dạng nửa khoảng là $[2; 6)$. Chọn \textbf{A}.
    \item \textbf{Câu 5:} Thay $(2; 1)$ vào hệ: $2(2) - 1 = 3 \ge 1$ và $2 + 1 = 3 \le 4$ (thỏa mãn cả hai bất phương trình). Chọn \textbf{C}.
    \item \textbf{Câu 6:} Thay $Q(1; 0)$ vào vế trái: $1 - 3(0) + 3 = 4 > 0$ (thỏa mãn). Chọn \textbf{D}.
    \item \textbf{Câu 7:} Thay $O(0; 0)$ vào vế trái: $2(0) + 0 - 4 = -4 \le 0$ (thỏa mãn). Miền nghiệm chứa gốc tọa độ $O(0; 0)$. Chọn \textbf{B}.
    \item \textbf{Câu 8:} Bất phương trình $3x - 2y + 6 \ge 0$ có đường thẳng bờ là phương trình dấu bằng: $3x - 2y + 6 = 0$. Chọn \textbf{D}.
    \item \textbf{Câu 9:} Với góc tù $\alpha \in (90^\circ; 180^\circ)$, ta có $\sin \alpha > 0$, $\cos \alpha < 0$, $\tan \alpha < 0$, $\cot \alpha < 0$. Chọn \textbf{A}.
    \item \textbf{Câu 10:} Áp dụng định lý côsin: $a^2 = b^2 + c^2 - 2bc \cos A = 6^2 + 8^2 - 2 \cdot 6 \cdot 8 \cdot \dfrac{1}{2} = 52 \implies a = 2\sqrt{13}$. Chọn \textbf{A}.
    \item \textbf{Câu 11:} Vì $\sin 50^\circ = \cos 40^\circ$ nên $\sin^2 40^\circ + \sin^2 50^\circ = \sin^2 40^\circ + \cos^2 40^\circ = 1$. Lại có $\tan 30^\circ \cdot \cot 30^\circ = 1$. Vậy $P = 1 + 1 = 2$. Chọn \textbf{C}.
    \item \textbf{Câu 12:} $\sin 150^\circ = \sin(180^\circ - 150^\circ) = \sin 30^\circ = \dfrac{1}{2}$. Chọn \textbf{D}.
\end{itemize}

\vspace{0.2cm}
\textbf{PHẦN II. Trắc nghiệm đúng/sai}
\begin{itemize}[leftmargin=0.6cm, itemsep=3pt]
    \item \textbf{Câu 13:}
    \begin{itemize}
        \item a) Thay tọa độ $M(1; 1)$ vào hệ: $1 \ge 0, 1 \ge 0, 1+1=2 \le 4, 1+2(1)=3 \le 6$ (thỏa mãn) $\Rightarrow$ \textbf{Đúng}.
        \item b) Cho $y = 0$, ta có $x \le 4$ và $x \le 6 \implies x \le 4$. Đỉnh trên trục $Ox$ là $A(4; 0)$ $\Rightarrow$ \textbf{Đúng}.
        \item c) Tọa độ giao điểm $B$ là nghiệm của hệ $\begin{cases} x + y = 4 \\ x + 2y = 6 \end{cases} \iff \begin{cases} x = 2 \\ y = 2 \end{cases}$. Điểm $B(2; 2)$ $\Rightarrow$ \textbf{Đúng}.
        \item d) Các đỉnh của tứ giác là $O(0; 0), A(4; 0), B(2; 2), C(0; 3)$. Tính giá trị $F(x, y) = 2x + 3y$:
        $F(O) = 0$; $F(A) = 8$; $F(B) = 2(2) + 3(2) = 10$; $F(C) = 2(0) + 3(3) = 9$.
        Giá trị lớn nhất là $10$ (tại $B(2; 2)$). Mệnh đề ghi $12$ là $\Rightarrow$ \textbf{Sai}.
    \end{itemize}

    \item \textbf{Câu 14:}
    \begin{itemize}
        \item a) $p = \dfrac{7 + 5 + 6}{2} = 9$ $\Rightarrow$ \textbf{Đúng}.
        \item b) $S = \sqrt{9(9 - 7)(9 - 5)(9 - 6)} = \sqrt{9 \cdot 2 \cdot 4 \cdot 3} = \sqrt{216} = 6\sqrt{6}$ $\Rightarrow$ \textbf{Đúng}.
        \item c) $r = \dfrac{S}{p} = \dfrac{6\sqrt{6}}{9} = \dfrac{2\sqrt{6}}{3}$ $\Rightarrow$ \textbf{Đúng}.
        \item d) $R = \dfrac{abc}{4S} = \dfrac{7 \cdot 5 \cdot 6}{4 \cdot 6\sqrt{6}} = \dfrac{35}{4\sqrt{6}} = \dfrac{35\sqrt{6}}{24}$. Mệnh đề ghi $R = \dfrac{35\sqrt{6}}{12}$ là $\Rightarrow$ \textbf{Sai}.
    \end{itemize}
\end{itemize}

\vspace{0.2cm}
\textbf{PHẦN III. Trắc nghiệm trả lời ngắn}
\begin{itemize}[leftmargin=0.6cm, itemsep=3pt]
    \item \textbf{Câu 15:} Gọi $T$ là tập hợp học sinh thích môn Toán, $V$ là tập hợp học sinh thích môn Ngữ văn.
    Số học sinh thích ít nhất một trong hai môn là: $n(T \cup V) = 40 - 5 = 35$.
    Áp dụng công thức cộng số phần tử:
    \[
    n(T \cup V) = n(T) + n(V) - n(T \cap V) \implies 35 = 25 + 20 - n(T \cap V) \implies n(T \cap V) = 10.
    \]
    Vậy có $10$ học sinh thích cả hai môn Toán và Ngữ văn.
    \item \textbf{Câu 16:} Miền nghiệm của hệ là miền tứ giác $OABC$ với các đỉnh:
    $O(0; 0)$, $A(3; 0)$ (giao của $3x + y = 9$ với $Ox$), $B(2; 3)$ (giao của hai đường thẳng $x + y = 5$ và $3x + y = 9$), $C(0; 5)$ (giao của $x + y = 5$ với $Oy$).
    Tính giá trị $F(x, y) = 4x + 3y$ tại các đỉnh:
    \begin{itemize}
        \item $F(O) = 0$;
        \item $F(A) = 4(3) + 3(0) = 12$;
        \item $F(B) = 4(2) + 3(3) = 17$;
        \item $F(C) = 4(0) + 3(5) = 15$.
    \end{itemize}
    Giá trị lớn nhất là $17$ (đạt tại đỉnh $B(2; 3)$).
    \item \textbf{Câu 17:} Vì $0^\circ < \alpha < 90^\circ$ nên $\cos \alpha > 0$. Ta có:
    \[
    \cos \alpha = \sqrt{1 - \sin^2 \alpha} = \sqrt{1 - \dfrac{25}{169}} = \dfrac{12}{13}, \quad \tan \alpha = \dfrac{\sin \alpha}{\cos \alpha} = \dfrac{5/13}{12/13} = \dfrac{5}{12}.
    \]
    Do đó: $T = 13 \cdot \dfrac{12}{13} + 12 \cdot \dfrac{5}{12} = 12 + 5 = 17$.
    \item \textbf{Câu 18:} Áp dụng định lý côsin trong tam giác $\triangle ABC$:
    \begin{align*}
    AB^2 &= CA^2 + CB^2 - 2 CA \cdot CB \cdot \cos 120^\circ \\
    &= 60^2 + 100^2 - 2(60)(100)\left(-\dfrac{1}{2}\right) = 3600 + 10000 + 6000 = 19600.
    \end{align*}
    Suy ra $AB = \sqrt{19600} = 140\text{ m}$.
\end{itemize}

\vspace{0.2cm}
\textbf{PHẦN IV. Tự luận}
\begin{itemize}[leftmargin=0.6cm, itemsep=4pt]
    \item \textbf{Câu 19:}
    Để tập hợp $A = [m - 3; m + 2]$ là tập con của tập hợp $B = [-1; 5)$, ta cần và đủ:
    \[
    \begin{cases}
    m - 3 \ge -1 \\
    m + 2 < 5
    \end{cases}
    \iff
    \begin{cases}
    m \ge 2 \\
    m < 3
    \end{cases}
    \iff 2 \le m < 3.
    \]
    Vậy tất cả các giá trị của tham số $m$ thỏa mãn là $m \in [2; 3)$.

    \item \textbf{Câu 20:}
    \begin{enumerate}
        \item Nửa chu vi tam giác $ABC$ là $p = \dfrac{13 + 14 + 15}{2} = 21$.
        Diện tích tam giác tính theo công thức Heron:
        \[
        S = \sqrt{21(21 - 13)(21 - 14)(21 - 15)} = \sqrt{21 \cdot 8 \cdot 7 \cdot 6} = \sqrt{7056} = 84.
        \]
        \item Mặt khác, $S = \dfrac{1}{2} a \cdot h_a \implies h_a = \dfrac{2S}{a} = \dfrac{2 \cdot 84}{13} = \dfrac{168}{13} \approx 12{,}9$.
    \end{enumerate}

    \item \textbf{Câu 21:}
    Trong tam giác $\triangle ABC$, góc ngoài tại đỉnh $B$ là $\widehat{CBD} = 55^\circ$, do đó:
    \[
    \widehat{ACB} = \widehat{CBD} - \widehat{CAD} = 55^\circ - 35^\circ = 20^\circ.
    \]
    Áp dụng định lý sin trong tam giác $\triangle ABC$:
    \[
    \dfrac{BC}{\sin \widehat{CAD}} = \dfrac{AB}{\sin \widehat{ACB}} \implies BC = \dfrac{AB \cdot \sin 35^\circ}{\sin 20^\circ} = \dfrac{30 \cdot \sin 35^\circ}{\sin 20^\circ} \approx 50{,}311\text{ m}.
    \]
    Trong tam giác vuông $\triangle BCD$ vuông tại $D$:
    \[
    h = CD = BC \cdot \sin 55^\circ \approx 50{,}311 \cdot \sin 55^\circ \approx 41{,}212\text{ m}.
    \]
    Làm tròn đến một chữ số thập phân, chiều cao cột phát sóng là $41{,}2\text{ m}$.

    \item \textbf{Câu 22:}
    Gọi $x, y$ lần lượt là số chiếc bàn và số chiếc ghế mà xưởng sản xuất trong một tuần ($x, y \in \mathbb{N}$).
    \begin{itemize}
        \item Lượng gỗ sử dụng: $2x + y \le 80$.
        \item Số giờ công lao động sử dụng: $3x + 2y \le 130$.
    \end{itemize}
    Ta có hệ bất phương trình ràng buộc:
    \[
    \begin{cases}
    x \ge 0,\; y \ge 0 \\
    2x + y \le 80 \\
    3x + 2y \le 130
    \end{cases}
    \]
    Miền nghiệm của hệ là miền tứ giác $OABC$ với các đỉnh:
    \begin{itemize}
        \item $O(0; 0)$;
        \item $A(40; 0)$ (giao điểm của $2x + y = 80$ với trục $Ox$);
        \item $B(30; 20)$ (nghiệm của hệ $\begin{cases} 2x + y = 80 \\ 3x + 2y = 130 \end{cases}$);
        \item $C(0; 65)$ (giao điểm của $3x + 2y = 130$ với trục $Oy$).
    \end{itemize}
    Tổng lợi nhuận thu được là $F(x, y) = 400x + 250y$ (nghìn đồng). Tính giá trị của $F$ tại các đỉnh:
    \begin{itemize}
        \item $F(O) = 0$;
        \item $F(A) = 400(40) = 16\,000$ (nghìn đồng) = $16$ triệu đồng;
        \item $F(B) = 400(30) + 250(20) = 12\,000 + 5\,000 = 17\,000$ (nghìn đồng) = $17$ triệu đồng;
        \item $F(C) = 250(65) = 16\,250$ (nghìn đồng) = $16{,}25$ triệu đồng.
    \end{itemize}
    Giá trị lớn nhất của $F(x, y)$ đạt được tại đỉnh $B(30; 20)$.\\
    Vậy xưởng cần sản xuất \textbf{30 chiếc bàn} và \textbf{20 chiếc ghế} để thu được lợi nhuận lớn nhất là \textbf{17.000.000 đồng}.
\end{itemize}

\end{document}
